\documentclass[10pt,a4paper]{article}
\usepackage[margin=18mm]{geometry}
\usepackage{amsmath,amssymb,booktabs,array}
\usepackage{xcolor}
\usepackage{cite}
\renewcommand{\citepunct}{], [}
\renewcommand{\citedash}{]--[}
\input{glyphtounicode}
\pdfgentounicode=1
\usepackage[hidelinks]{hyperref}
\hypersetup{pdftitle={An elementary VSL lane-drop task: model and sensitivity},
 pdfauthor={Differentiable traffic task review},pdfsubject={Compact methodology}}
\setlength{\parindent}{0pt}
\setlength{\parskip}{3pt}
\setlength{\abovedisplayskip}{4pt}
\setlength{\belowdisplayskip}{4pt}
\setlength{\abovedisplayshortskip}{3pt}
\setlength{\belowdisplayshortskip}{3pt}
\newcommand{\pos}[1]{[#1]_+}
\newcommand{\bud}{\mathcal B}
\newcommand{\dd}{\,\mathrm d}
\newcommand{\head}[1]{\par\vspace{4pt}\noindent\textbf{#1}\par\nobreak\vspace{1pt}}
\newcommand{\rev}[1]{\textcolor{red}{#1}}
\newenvironment{points}{\begin{list}{$\bullet$}{%
 \setlength{\leftmargin}{1.4em}\setlength{\labelwidth}{0.7em}%
 \setlength{\labelsep}{0.4em}\setlength{\itemsep}{2pt}%
 \setlength{\parsep}{0pt}\setlength{\topsep}{3pt}}}{\end{list}}
\begin{document}
\begin{center}
{\large\bfseries An elementary VSL lane-drop task: model and sensitivity}
\end{center}
\vspace{-5pt}
\rev{\textit{Revisions responding to feedback 11 are shown in red.}}

\head{1. Research question}
\begingroup\color{red}
This elementary variable-speed-limit (VSL) task separates three questions:
\begin{points}
\item \textbf{Information:} Which omitted lane-loading information changes
merge service when ending-lane traffic and total continuing traffic agree?
\item \textbf{Decision:} Can retaining that information change the useful
common-speed advice after both travel and all waiting time are counted?
\item \textbf{Gradient:} Does the complete travel-time sensitivity predict
beneficial speed changes, checked against actual objective changes and
equal-information optimizer references?
\end{points}
\endgroup

\head{2. Task and global assumptions}
\rev{\textbf{Concrete task.} Control one permanent three-to-two lane drop
for one specified finite cohort, using one constant common speed advice.}
\rev{The road has} ending lane E, adjacent continuing lane A
and far continuing lane B. Preparation ends at $L_p$; a terminal of length
$b$ ends at the choke $L=L_p+b$. At time zero, independent lane densities
$(e_0,a_0,b_0)\geq0$ are uniform on the common footprint $[0,\ell]$,
$0<\ell\leq L_p$. Set $r=e_0+a_0$ and total mass
$M=\ell(e_0+a_0+b_0)>0$. The terminal and recovery history start empty;
there are no later external arrivals. Choose one constant common advice
$v\in[v_-,v_+]$, $0<v_-<v_+$, with fixed initial state, geometry and physical parameters.
The objective counts every initial unit's time until it crosses the choke.

\begingroup\color{red}
\textbf{Global modeling assumptions.} These specify the traffic represented
by this task, separately from the sufficient operating bounds below.
\begin{points}
\item \textbf{MA1: continuous, homogeneous traffic.} Each lane carries
continuous vehicle-equivalent mass. The remaining preparation body retains
lane densities without internal longitudinal variation or individual maneuvers.
\item \textbf{MA2: advice acts through traffic state.} Unblocked traffic
follows the common advice. Backward-wave speed $w$, jam density $k_j$,
geometry and physical interaction/recovery coefficients stay fixed as advice
varies; merge service responds to the resulting local traffic state.
\item \textbf{MA3: two-stage lane coupling.} E can transfer into A before
the terminal; successful early transfer becomes ordinary A. Residual E
traverses the terminal without ordinary lane changes and merges at the choke.
B has separate service; no other lateral exchange is represented.
\item \textbf{MA4: point-service closure.} Current A/E conditions determine
priority/cooperation, maneuver work and representative boundary speeds.
Those speeds describe free or queued service, rather than individual trajectories.
\item \textbf{MA5: recovery memory.} Only forced E discharge generates
recovery loss. One shared A/E budget depends on the retained ages and
insertion speeds of the most recent fixed amount of E mass.
\end{points}
\endgroup

Triangular-diagram and cumulative-boundary ingredients follow the link
transmission model (LTM)~\cite{jin}; the advice-response family is selected
here, with algebraic capacity/inverse maps:
\begin{equation}
 K=wk_j,\quad k_c(v)=\frac{K}{w+v},\quad C(v)=vk_c(v),\qquad
 c(u)=\frac{Ku}{w+u},\quad U(g)=\frac{wg}{K-g}.
 \label{eq:fd}
\end{equation}
\rev{Here $k_c$ is critical density, $C$ capacity, $c(u)$ the
congested-branch flow associated with speed $u$, and $U$ its inverse.}
$R$ and $R_B$ are undisturbed A/E and B service ceilings, respectively.
Use positive physical scales and the sufficient operating domain
\begin{equation}
 e_0,a_0,b_0\leq k_c(v_+),\qquad
 v_+r<R\leq C(v_-),\qquad v_+b_0\leq R_B\leq C(v_-),\qquad M<k_jb.
 \label{eq:domain}
\end{equation}
It keeps preparation free, terminal receiving available and B unqueued.
Work and recovery choices are specified below; $[x]_+=\max(x,0)$.

\begingroup\color{red}
\textbf{Why this test.} Lane-based conservation and finite receiving
constraints~\cite{subraveti,carey}, together with insertion-speed effects
on discharge~\cite{leclercq2011,leclercq}, motivate conserved, feasible
traffic and state-sensitive merge service. Strong aggregate VSL models
already retain traffic history and capacity drop~\cite{wei}, but do not
explicitly resolve the A/B split. Lane-resolved models~\cite{subraveti,carey}
supply that state; their traffic predictions alone do not answer this task's
matched-information control/gradient questions. We therefore require
complete travel-time accounting through clearance, lawful history inference,
and separate checks of information value and equal-information gradient
usefulness, including a sufficient compact peer and scalar search.
\endgroup

\newpage
\head{3. Traffic model}
\begingroup\color{red}
Picture E traffic approaching the lane end. Some finds room in A early and
continues as ordinary A traffic. What remains in E crosses a short terminal
and must merge at the choke, competing with A for passage. A slow forced
entry can keep restricting following traffic even after it has crossed the
completion line; the restriction fades as downstream traffic accelerates.
The four components below represent these early opportunities, travel to
the merging point, competition for passage, and the lingering recovery effect.
\endgroup

\textbf{3.1. Pre-choke merging.}
Let $e(t),a(t)$ be the remaining preparation body's E/A densities.
Own A flow uses receiving first. E transfer is limited by its donor offer
and by remaining adjacent service.
\rev{Paired conservative sources follow~\cite{subraveti}; receiving work
and own-lane priority draw on~\cite{carey}. The homogeneous minimum and
interaction length below are defined here:}
\begin{equation}
 S(v,a)=\min\{C(v),w(k_j-a)\},\qquad
 f=\min\!\left\{ve,\frac{S(v,a)-va}{\gamma}\right\},\qquad
 \dot e=-f/h,\quad \dot a=f/h.
 \label{eq:prep}
\end{equation}
Here $h>0$ is a physical interaction length, $\gamma\geq1$ is a receiving-work
factor, and $f/h$ is the lateral density source. Thus $e+a=r$ and $b_0$
stays constant. Donor-limited decay is exponential per traveled distance;
advice affects preparation through normalized receiving headroom $k_c(v)$.

\textbf{3.2. Finite terminal and the merging point.}
\rev{The terminal separates reaching the end of preparation from being
ready to pass the choke: traffic first travels distance $b$, then waits
only if point service cannot accept it.}
For $j\in\{E,A,B\}$, let $F_j,P_j,Y_j$ be cumulative terminal entry,
service-eligible mass and departure.
\rev{The initially empty count histories represent only this cohort.
Boundary travel follows LTM~\cite{jin}; the preparation handoff is defined
here. The cohort enters during $(L_p-\ell)/v<t<L_p/v$:}
\begin{equation}
 \dot F_E=ve(t),\quad \dot F_A=v[r-e(t)],\quad \dot F_B=vb_0;
 \qquad P_j(t)=F_j(t-b/v).
 \label{eq:arrivals}
\end{equation}
Define physical terminal stock $N_j=F_j-Y_j$, eligible waiting mass
$q_j=P_j-Y_j$, offer $\lambda_j=\dot P_j$ and discharge $g_j=\dot Y_j$.
Departures satisfy $0\leq Y_j\leq P_j\leq F_j$, with
$\dot q_j=\lambda_j-g_j$. Domain~\eqref{eq:domain} protects finite storage
and gives $Y_B=P_B$. Early-merged E becomes ordinary A; residual E travels
through the terminal and merges at the choke.

\newpage
\textbf{3.3. Forced merging and service sharing.}
\begingroup\color{red}
The point has a shared work budget $\bud$, measured as A-equivalent service
per unit time. A costs one work unit per mass; E's maneuver costs more
when A is faster. For lane $j=A,E$, $u_j$ is the model's representative
speed at the merging boundary: $u_j=v$ for unqueued traffic and
$u_j=U(g_j)$ for queued traffic. The queued relation is an equilibrium
closure, not an individual speed trajectory.
Receiving-work accounting is adapted from~\cite{carey}; the specific
speed-mismatch work and allocation below are defined here:
\endgroup
\begin{equation}
 \begin{gathered}
 c_E=c_0+\chi g_A\pos{u_A-u_E},\qquad
 g_A+c_Eg_E\leq\bud,\\
 \rev{H(g;g_A,u_A)=g\cdot\bigl(c_0+\chi g_A\pos{u_A-U(g)}\bigr)=x.}
 \end{gathered}
 \label{eq:work}
\end{equation}
$c_0\geq1$ is baseline E work and $\beta\geq0$ sets mismatch strength.
Set $\chi=\beta/[C(v_+)v_+]$ once and hold it fixed.
Require $c_0>\beta(1+v_+/w)$ so each feasible root $H(g)=x$
is unique.
\rev{When E is queued, \textbf{solve for $g=g_E$}: A's flow $g_A$, speed
$u_A$ and allocated E work $x$ come from the sharing table. $H$ is the fully
specified work consumed by a candidate E flow $g$; the dot denotes ordinary
multiplication of that flow by its work per unit mass.}
A fitting free A offer has
priority; queued A cooperates through equal work and unused-work redistribution:
\begin{center}
\setlength{\tabcolsep}{5pt}
\begin{tabular}{@{}ll@{}}
\toprule
A, E status & Selected service and free-admission condition\\
\midrule
Free, free & $g_A=\lambda_A,\ g_E=\lambda_E$ if
$\lambda_A+c_0\lambda_E\leq\bud$.\\
Free, queued & $g_A=\lambda_A\leq\bud,\quad
H(g_E;\lambda_A,v)=\bud-\lambda_A$.\\
Queued, queued & $g_A=\bud/2,\quad
H(g_E;\bud/2,U(\bud/2))=\bud/2$.\\
Queued, free & $g_E=\lambda_E,\ g_A=\bud-c_0\lambda_E$ if
$c_0\lambda_E\leq\bud/2$.\\
\bottomrule
\end{tabular}
\end{center}
\begingroup\color{red}
A queue starts only when arriving traffic strictly exceeds feasible service.
When all waiting traffic has passed, the lane becomes unqueued if its
current arrivals can pass immediately; equality permits that continuation.
After an arrival-rate change, service follows the new arrivals immediately.
Recovery is a separate physical aftereffect: clearing a queue does not erase
the influence of traffic already discharged. These transition rules and
the sharing table are defined here.
\endgroup

\textbf{3.4. Recovery and the shared budget.}
\begingroup\color{red}
Crossing the choke completes an E entry's own trip, but a slow insertion
can still hold back the A/E traffic waiting behind it. Downstream acceleration
gradually releases that restriction. Point-merge studies supply this
mechanism~\cite{leclercq2011,leclercq}. For one entry with insertion speed
$u$, let $\theta\geq0$ be time since crossing and $a_r>0$ the recovery
acceleration. Differentiating the isolated cumulative passing expression
in~\cite{leclercq} gives $c(V)$; capping against undisturbed service $R$
defines the loss below:
\endgroup
\begin{equation}
 V(\theta;u)=\sqrt{(w+u)^2+2wa_r\theta}-w,\qquad
 D(\theta;u)=\pos{R-c(V(\theta;u))}.
 \label{eq:profile}
\end{equation}
\begingroup\color{red}
$V$ summarizes how the downstream acceleration and backward-wave propagation
restore service at the choke. It is an equivalent speed in the boundary-flow
map $c$, not the discharged vehicle's actual speed. $D$ is the amount of
the undisturbed A/E service $R$ still unavailable because of this insertion;
it vanishes once $c(V)$ reaches $R$.
To combine entries, $\xi$ labels cumulative discharged E mass,
$T(\xi)$ is that mass's choke-crossing time, and $u_E(T(\xi))$ its insertion
speed. Set $m=1$ vehicle-equivalent. Our recent-mass rule averages the
remaining losses over at most the latest $m$ E mass and subtracts the
reading $p_{\rm rec}$ from $R$ to obtain the shared budget:
\endgroup
\begin{equation}
 p_{\rm rec}(t)=\frac1m\int_{\pos{Y_E(t)-m}}^{Y_E(t)}
 D\bigl(t-T(\xi);u_E(T(\xi))\bigr)\dd\xi,\qquad
 \bud(t)=\pos{R-p_{\rm rec}(t)}.
 \label{eq:recovery}
\end{equation}
\rev{All retained profiles recover with age. Once $m$ has passed, newer
E replaces the oldest retained mass; the denominator remains $m$ even before
fill. A consumes the shared budget but adds no separate loss profile.
The capped single-entry profile is derived/adapted from~\cite{leclercq};
the E-only source and recent-mass renewal are defined here.}

\head{4. How the task tests information}
Compare the same initial E mass and continuing total with different A/B
splits. A busier A can restrict preparation, leaving more E for forced work
and recovery. Advice trades these effects against ordinary travel time.
With a fixed horizon $H_*$ beyond cohort clearance, minimize complete
\rev{total system time (TST)}:
\begin{equation}
 J(v)=\int_0^{H_*}\!\left[M-\sum_jY_j(t)\right]\dd t
 =\frac{M(L-\ell/2)}v+\int_0^{H_*}(q_A+q_E)\dd t.
 \label{eq:objective}
\end{equation}
The predictors share geometry and physical parameters and may use all
inference from admitted predecision histories: aggregate density, flow,
speed and cumulative counts, including any permitted E-resolved history.
One knows initial E and combined A+B mass; the other knows each lane.
An equal-information compact peer adds just initial A mass to the former,
reconstructing all three initial densities and the same dynamics.
\rev{A separate optimizer comparison uses the same information and model
for gradient-based advice and bounded scalar search.}

\newpage
\head{5. Complete sensitivity}
The model-specific sensitivities are derived here at fixed initial traffic
and physical parameters.
\rev{Advice changes preparation and arrival timing; these change passage,
insertion speeds, recovery and waiting. Differentiate that whole path,
including when queues form/clear and old recovery entries leave the account.}
\rev{Here $H$ is the scalar service-work function defined in Section~3.3,
Eq.~\eqref{eq:work}. The terms $H_g$, $H_{g_A}$ and $H_{u_A}$ are its first
partial derivatives with respect to E discharge, A discharge and A's
representative boundary speed, respectively, evaluated at the current
service solution. Each measures how required work changes with one argument
while the others are held fixed.}
A subscript $v$ denotes advice sensitivity. On smooth branches,
\begin{equation}
 g_{E,v}=\frac{x_v-H_{g_A}g_{A,v}-H_{u_A}u_{A,v}}{H_g},\qquad H_g>0.
 \label{eq:rootgradient}
\end{equation}
\begingroup\color{red}
Let $s_m$ be the crossing time of the oldest retained E mass. For a filled
window with $g_E(s_m)>0$, $Y_E(s_m)=Y_E(t)-m$; write $Z_E=Y_{E,v}$.
The cutoff relation below is derived here. Differentiate~\eqref{eq:recovery}
including insertion speeds, flows, retirement and source events. A branch
switch at $t=\tau$ satisfies a condition $G(t,v)=0$ along the predicted
traffic/history. For an isolated crossing with $G_t^-\ne0$, the remaining
event corrections follow~\cite{kong}. Here $X$ collects continuous traffic
coordinates and $-/+$ mean just before/after the switch:
\endgroup
\begin{equation}
 \begin{gathered}
 g_E(s_m)s_{m,v}=Z_E(t)-Z_E(s_m),\qquad \tau_v=-G_v^-/G_t^-,\\
 X_v^+=X_v^-+(\dot X^- -\dot X^+)\tau_v.
 \end{gathered}
 \label{eq:eventgradient}
\end{equation}
Thus the complete derivative, on a fixed regular event/history sequence, is
\begin{equation}
 J'(v)=-\int_0^{H_*}\!\sum_jY_{j,v}(t)\dd t
 =-\frac{M(L-\ell/2)}{v^2}+\int_0^{H_*}(q_{A,v}+q_{E,v})\dd t.
 \label{eq:objectivegradient}
\end{equation}
\makeatletter
\renewcommand{\section}{\@ifstar{\head}{\head}}
\makeatother
% Main methodology target: at most 3.5 pages; references are excluded.
\begin{thebibliography}{7}
\setlength{\itemsep}{0pt}\setlength{\parskip}{0pt}
\bibitem{jin} W.-L. Jin, ``Continuous formulations and analytical properties
of the link transmission model,'' \emph{Transp. Res. B, Methodol.}, vol.~74,
pp.~88--103, 2015, doi: \href{https://doi.org/10.1016/j.trb.2014.12.006}{10.1016/j.trb.2014.12.006}.
\bibitem{subraveti} H.~H.~S. Nagalur Subraveti, V.~L. Knoop, and B. van Arem,
``First order multi-lane traffic flow model---an incentive based macroscopic
model to represent lane change dynamics,'' \emph{Transportmetrica B, Transp.
Dyn.}, vol.~7, no.~1, pp.~1758--1779, 2019, doi:
\href{https://doi.org/10.1080/21680566.2019.1700846}{10.1080/21680566.2019.1700846}.
\bibitem{carey} M. Carey, C. Balijepalli, and D. Watling, ``Extending the cell
transmission model to multiple lanes and lane-changing,'' \emph{Netw. Spat.
Econ.}, vol.~15, pp.~507--535, 2015, doi:
\href{https://doi.org/10.1007/s11067-013-9193-7}{10.1007/s11067-013-9193-7}.
\bibitem{leclercq2011} L. Leclercq, J.~A. Laval, and N. Chiabaut,
``Capacity drops at merges: An endogenous model,'' \emph{Procedia Soc.
Behav. Sci.}, vol.~17, pp.~12--26, 2011, doi:
\href{https://doi.org/10.1016/j.sbspro.2011.04.505}{10.1016/j.sbspro.2011.04.505}.
\bibitem{leclercq} L. Leclercq, V.~L. Knoop, F. Marczak, and S.~P. Hoogendoorn,
``Capacity drops at merges: New analytical investigations,'' \emph{Transp.
Res. C, Emerg. Technol.}, vol.~62, pp.~171--181, 2016, doi:
\href{https://doi.org/10.1016/j.trc.2015.06.025}{10.1016/j.trc.2015.06.025}.
\bibitem{wei} \rev{L. Wei, Y. Han, and M. Wang, ``A first-order link-based flow
model with variable speed limits and capacity drops for freeway networks,''
\emph{Comput.-Aided Civ. Infrastruct. Eng.}, vol.~40, no.~21,
pp.~3200--3217, 2025, doi:
\href{https://doi.org/10.1111/mice.13492}{10.1111/mice.13492}.}
\bibitem{kong} N.~J. Kong, J.~J. Payne, J. Zhu, and A.~M. Johnson,
``Saltation matrices: The essential tool for linearizing hybrid dynamical
systems,'' \emph{Proc. IEEE}, vol.~112, no.~6, pp.~585--608, Jun. 2024, doi:
\href{https://doi.org/10.1109/JPROC.2024.3440211}{10.1109/JPROC.2024.3440211}.
\end{thebibliography}
\end{document}
